% Generated draft: requires mathematical and source review.
\documentclass[11pt]{article}
\usepackage[margin=25mm]{geometry}
\usepackage{fontspec}
\setmainfont{Latin Modern Roman}
\usepackage{amsmath,amssymb,mathtools,unicode-math}
\setmathfont{Latin Modern Math}
\usepackage{enumitem,needspace,xurl,hyperref}
\setlength{\parindent}{0pt}
\setlength{\parskip}{5pt}
\setlength{\emergencystretch}{3em}
\newcommand{\N}{\mathbb{N}}
\newcommand{\Z}{\mathbb{Z}}
\newcommand{\Q}{\mathbb{Q}}
\newcommand{\R}{\mathbb{R}}
\newcommand{\C}{\mathbb{C}}
\newcommand{\F}{\mathbb{F}}
\newcommand{\A}{\mathbb{A}}
\newcommand{\PP}{\mathbb{P}}
\newcommand{\E}{\mathbb{E}}
\newcommand{\eps}{\varepsilon}
\newcommand{\dd}{\,\mathrm{d}}
\begin{document}
\section*{UnsolvedMath ID: 30001965}
\textbf{Source problem code:} catalogue entry
\subsection{Definitions and mathematical statement}
The supplied record poses the following historical and methodological question:
\begin{quote}
Can historical events in the development of mathematics be understood more plausibly by treating tacit knowledge as a form of knowledge?
\end{quote}
The record supplies no mathematical definition of tacit knowledge, no specified collection of historical events, no competing interpretations, and no criterion for comparing their plausibility. It also specifies no quantifier over historical events. Consequently, a precise mathematical proposition cannot be recovered from the supplied evidence without adding substantive assumptions. The question is preserved as a discussion prompt; no mathematical formalization is asserted.

The supplied statement-verification note characterizes the workshop discussion as historical, pedagogical, or methodological rather than as an unresolved mathematical proposition. The record has statement status \texttt{unrecoverable}, supplies no clean statement, and retains the original question for provenance. Its supplied \texttt{open} label is a registry classification, not an independently verified mathematical status.
\subsection{Short English statement}
The question asks whether recognizing tacit knowledge as knowledge can make interpretations of mathematical history more plausible. The supplied record does not define the concepts or comparison criteria needed to turn this discussion prompt into a precise mathematical statement.
\subsection{Sources}
\begin{itemize}
\item Local UnsolvedMath registry, permanent ID \texttt{30001965}, title \emph{Tacit Knowledge in Mathematical History}, source code \texttt{catalogue entry}, in the \emph{Oberwolfach Reports Open Problems} collection. The supplied record assigns proposed year 2012 but identifies no individual proposer.
\item Supplied source citation: \emph{Explicit Versus Tacit Knowledge in Mathematics} (2012), \emph{Oberwolfach Reports}, pp.~131--210, DOI \texttt{10.4171/owr/2012/04}. Supplied source URL and literature-source link, labeled ``original Oberwolfach report'': \url{https://doi.org/10.4171/owr/2012/04}. The supplied background describes these reports as documenting mathematical workshops and their open-problem discussions.
\item Literature assessment dated exactly \texttt{2026-08-21}: the supplied curation labels the record \texttt{open}, while explaining that the extracted question concerns historical interpretation of tacit knowledge and is not a claim that can be marked solved or disproved. This is supplied dated curation, not independently verified. The background calls for independent verification of time-sensitive status and cited claims.
\item The supplied statement-verification note describes the official workshop report as treating this item as a historical, pedagogical, or methodological discussion prompt, with no separate clean mathematical statement promoted. This characterization is preserved as supplied source-verification information, not as an independent inspection of the report.
\end{itemize}
\end{document}
